\section{Tiên đề tách}

\

Giả sử $x$ là một tập hợp và $\cauF()$ có nghĩa với toàn bộ các phần tử của $x$. Khi đó, tiên đề tách khẳng định rằng tồn tại một tập hợp $y$ bao gồm tất cả các phần tử $z \in x$ thỏa mãn $\cauF(z)$ và chỉ những phần tử đó. Nói cách khác,
$$
    \defMath{\forall x, \exists y, \forall z, \left( z \in y \iff z \in x \land \textbf{\cauF}(z) \right)}.
$$

\begin{figure}[H]
    \begin{center}
    \begin{tikzpicture}[>=stealth]

    \draw[thick] (-3,0) ellipse (2.2cm and 2.8cm);
    \node at (-3, 3.2) {Tập hợp $x$};

    \node[circle, fill=colorEmphasisCyan, inner sep=\pointSize, label=left:$z_1$] (z1_left) at (-3.7, 0.8) {};
    \node[circle, fill=colorEmphasisCyan, inner sep=\pointSize, label=left:$z_3$] (z3_left) at (-2.5, -0.5) {};

    \node[circle, inner sep=\pointSize, label=left:{$z_2$}] (z2_left) at (-2, 0.8) {};
    \node[circle, inner sep=\pointSize, label=left:{$z_4$}] (z4_left) at (-3.8, -1.2) {};

    \draw[dashed, colorEmphasisCyan, thin] plot [smooth cycle, tension=1] coordinates {
        (-4.6, 0.5)
        (-3.5, 1.8)
        (-3, 0.4)
        (-2, -0.2)
        (-2.5, -1.2) 
    };
    \node at (-3.1, 2.1) {$\textcolor{colorEmphasisCyan}{\cauF(}z\textcolor{colorEmphasisCyan}{)}$};


    \draw[->, ultra thick, colorEmphasis] (0,0) -- (2,0) 
        node[midway, above, align=center] {Tách bởi\\ $\textcolor{colorEmphasisCyan}{\cauF(}z\textcolor{colorEmphasisCyan}{)}$};


    \draw[thick] (4.5,0) ellipse (1.5cm and 2.2cm);
    \node at (4.5, 2.6) {$\{z \in x \mid \textcolor{colorEmphasisCyan}{\cauF(}z\textcolor{colorEmphasisCyan}{)}\}$};

    % Chỉ các phần tử thỏa mãn mới xuất hiện ở đây
    \node[circle, fill=colorEmphasisCyan, inner sep=1.5pt, label=right:$z_1$] (z1_right) at (4.2, 0.8) {};
    \node[circle, fill=colorEmphasisCyan, inner sep=1.5pt, label=right:$z_3$] (z3_right) at (4.8, -0.5) {};


    \draw[->, thick, colorEmphasisCyan, shorten >= 5pt, shorten <= 5pt, dashed] 
        (z1_left) to[bend left=20] (z1_right);

    \draw[->, thick, colorEmphasisCyan, shorten >= 5pt, shorten <= 5pt, dashed] 
        (z3_left) to[bend right=15] (z3_right);

    \end{tikzpicture}
    \end{center}
    \caption{Minh họa tiên đề tách (các phần tử của $x$ thỏa mãn tính chất $\textcolor{colorEmphasisCyan}{\cauF(}z\textcolor{colorEmphasisCyan}{)}$ được tách ra).}
\end{figure}

Tiên đề ngoại diên khẳng định tập hợp $y$ là duy nhất. Thực vậy, giả sử tồn tại hai tập hợp $y_1$ và $y_2$ thỏa mãn điều kiện của tiên đề tách với cùng tập hợp $x$ và vị từ $\cauF()$. Xét một $z$ bất kì. Khi đó, định nghĩa theo tiên đề cho chúng ta
\begin{equation*}
    \begin{cases}
        z \in y_1 \iff z \in x \land \cauF(z) \\
        z \in y_2 \iff z \in x \land \cauF(z)
    \end{cases}.
\end{equation*}
Theo tính bắc cầu của phép đẳng giá, chúng ta thu được
\begin{equation*}
    \forall z, \left( z \in y_1 \iff z \in y_2 \right).
\end{equation*}
Theo tiên đề ngoại diên, $y_1 = y_2$. Gọi tập hợp đó là $y$. Một hệ quả khá hiển nhiên chính là $y \subseteq x$, bởi vì rõ ràng rằng, nếu $z \in y$, theo tiên đề tách, chúng ta phải có $z \in x$.

Trước khi đi vào những vấn đề khác, chúng ta sẽ trình bày về một cách viết tắt các mệnh đề trong giải tích vị từ bậc nhất. Bạn đọc có thể để ý, trong định nghĩa, chúng ta có một mệnh đề dạng
$$
   \forall z, \left(z \in x \implies V(z)\right).
$$
Dưới dạng ngôn ngữ, chúng ta sẽ nói:
\begin{center}
   ``Với mọi, nếu $z$ thuộc $x$ thì $V(z)$.'',
\end{center}
nhưng một cách phát biểu tự nhiên hơn sẽ là:
\begin{center}
   ``Với mọi $z$ thuộc $x$, $V(z)$.''
\end{center}
Từ cách nói đó, chúng ta có cách viết rút gọn với mọi tập hợp $x$ và vị từ $V()$ như sau:
\begin{equation*}
   \defMath{\forall z \in x, V(z)} \iff \forall z, \left(z \in x \implies V(z)\right).
\end{equation*}
Tương tự, với lượng từ tồn tại, chúng ta cũng có cách viết tắt:
\begin{equation*}
   \defMath{\exists z \in x, V(z)} \iff \exists z, \left(z \in x \land V(z)\right).
\end{equation*}

Quay trở lại với tập hợp $y$ trong tiên đề tách, chúng ta không thể kí hiệu bằng việc thêm một kí tự duy nhất như trước được nữa, mà do mỗi tập $y$ đều cần được xác định từ một vị từ $\cauF$, vì thế, chúng ta sẽ kí hiệu tập hợp $y$ như sau:
\begin{equation*}
   y = \defMath{\{ z \in x \mid \cauF(z) \}} = \defMath{\{ z \in x : \cauF(z) \}}.
\end{equation*}
và hệ quả sẽ được kí hiệu là
\begin{equation}
   \{ z \in x \mid \cauF(z) \} \subseteq x.
   \label{eq:li_thuyet_tap_hop:tien_de:tien_de_tach:he_qua_tap_con}
\end{equation}

Liệu bạn đọc còn nhớ về nghịch lí Rất-xen, liên quan đến một tập bao gồm tất cả các tập hợp không chứa chính nó. Có một cách đơn giản để đi vòng qua vấn đề này. Đó chính là không cho nó tồn tại. Để làm như vậy, chúng ta sẽ chứng minh không tồn tại tập hợp $\Omega$ sao cho $\forall x, x \in \Omega$. Thực vậy, giả sử tồn tại tập hợp $\Omega$ như vậy. Khi đó, theo tiên đề tách, tồn tại tập hợp $y = \{x \in \Omega \mid x \notin x \}$. Đặt câu hỏi, $y$ có thuộc $y$ không?
\begin{itemize}
    \item Nếu $y \in y$, thì theo định nghĩa của $y$, chúng ta có $y \notin y$, do $y$ chỉ chứa các tập không chứa chính nó, một mâu thuẫn.
    \item Nếu $y \notin y$, thì tương tự, theo định nghĩa của $y$, chúng ta có $y \in y$, một mâu thuẫn.
\end{itemize}
Nhìn chung, cả hai trường hợp đều đưa ra một điều vô lí. Sử dụng chứng minh bằng phản chứng, chúng ta có thể kết luận rằng không tồn tại tập hợp $\Omega$ sao cho $x \in \Omega$ với mọi $x$.

Giờ là lúc để làm một số ví dụ. Giả sử $x$ và $y$ là hai tập hợp. Nhận thấy rằng, ``$(z) \in y$'' là một vị từ theo hạng thức $z$. Áp dụng tiên đề tách, chúng ta có tập
\begin{equation*}
    \defMath{\{ z \in x \mid z \in y \}}.
\end{equation*}
Đây được gọi là \defText{giao} của $x$ và $y$, kí hiệu $\defMath{x \cap y}$. Nếu $x \cap y = \varnothing$, thì $x$ và $y$ được gọi là \defText{rời nhau}. Ngoài ra, có $z \in y$, thì chúng ta cũng có thể có $z \notin y$, điều đó tự nhiên dẫn đến \defText{hiệu} giữa $x$ và $y$ như sau:
\begin{equation*}
    \defMath{x \setminus y = x - y = \{ z \in x \mid z \notin y \}}.
\end{equation*}
$x \setminus y$ còn được gọi là \defText{phần bù tương đối} của $y$ trong $x$. Thường xuyên, chúng ta sẽ chỉ quan tâm các tính chất giới hạn trong một tập cụ thể, hiểu theo nghĩa, chúng ta sẽ chỉ quan tâm đến tập con của một tập hợp $x$ nào đó. Với giả thiết rằng $y \subseteq x$, chúng ta định nghĩa và kí hiệu \defText{phần bù} của $y$ là $$\defMath{\complement_x(y) = \complement(y) = y^c = \overline{y} = x \setminus y}.$$

Định nghĩa giao có thể được mở rộng lên nhóm nhiều tập hợp. Tiên đề hợp cho phép xây dựng tập $\bigcup x$, và ``$\forall a \in (x), z \in a$'' là một vị từ theo hạng thức $x$. Sử dụng tiên đề tách, chúng ta có tập
{
    \boldmath 
    $$\bigcap x = \bigcap_{a\in x} a = \left\{z \in \bigcup x \;\middle|\; \forall a \in x, z \in a \right\}$$
}%
và gọi nó là giao của các phần tử trong $x$. Bạn đọc cần phải chú ý một điều, để không xảy ra mâu thuẫn, định nghĩa mới phải tương thích với định nghĩa cũ. May mắn thay, $\bigcap \{x; y\} = x \cap y$.

Một phép toán nữa, khá ít dùng, được gọi là \defText{hiệu đối xứng}. Với hai tập hợp $x$ và $y$, hiệu đối xứng của chúng được định nghĩa là
\begin{equation*}
    \defMath{x \vartriangle y = (x \setminus y) \cup (y \setminus x)}.
\end{equation*}

Để kết thúc phần này, tác giả sẽ giới thiệu về cách biểu diễn tập hợp bằng sơ đồ Ven\footnote{John Venn (1834 -- 1923).}. Với hai tập hợp $x$ và $y$, chúng ta có thể biểu diễn chúng bằng hai hình chồng lên nhau, trong đó kết quả của các phép tính trên chúng được tô màu, với một sự minh hoạ cụ thể ở hình \ref{fig:li_thuyet_tap_hop:tien_de:tien_de_tach:ven_phep_tinh}.

\begin{figure}[H]
    \centering
    \begin{tikzpicture}
        \begin{scope}[shift={(0,0)}]
            % Tô màu cả 2 hình tròn
            \fill[color = colorEmphasis] (0,0) circle (0.7);
            \fill[color = colorEmphasis] (0.8,0) circle (0.7);
            
            % % Vẽ đường viền
            \draw (0,0) circle (0.7);
            \draw (0.8,0) circle (0.7);
            
            % Nhãn tên tập hợp (màu trắng)
            \node[above left] at (-0.4, 0.4) {$x$};
            \node[above right] at (1.2, 0.4) {$y$};
            
            % Nhãn phép toán
            \node[below] at (0.4, -0.9) {$x \cup y$};
        \end{scope}

        \begin{scope}[shift={(4,0)}]
            \begin{scope}
                \clip (0,0) circle (0.7);
                \fill[color = colorEmphasis] (0.8,0) circle (0.7);
            \end{scope}
            
            % Vẽ đường viền
            \draw (0,0) circle (0.7);
            \draw (0.8,0) circle (0.7);
            
            % Nhãn
            \node[above left] at (-0.4, 0.4) {$x$};
            \node[above right] at (1.2, 0.4) {$y$};
            \node[below] at (0.4, -0.9) {$x \cap y$};
        \end{scope}

        \begin{scope}[shift={(8, 0)}]
            % Tô màu tập x nhưng bỏ đi phần giao với y
            \begin{scope}
                \clip (0,0) circle (0.7);
                \fill[white] (0.8,0) circle (0.7); % Xóa phần giao bằng nền
            \end{scope}
            % Phủ lại pattern lên phần còn lại của x
            \begin{scope}
                \clip (0,0) circle (0.7);
                \draw (0.8,0) circle (0.7); 
                % Thủ thuật đảo ngược để tô phần x \ y
                \fill[color = colorEmphasis, even odd rule] (0,0) circle (0.7) (0.8,0) circle (0.7);
            \end{scope}
            
            % Vẽ đường viền
            \draw (0,0) circle (0.7);
            \draw (0.8,0) circle (0.7);
            
            % Nhãn
            \node[above left] at (-0.4, 0.4) {$x$};
            \node[above right] at (1.2, 0.4) {$y$};
            \node[below] at (0.4, -0.9) {$x \setminus y$};
        \end{scope}

        \begin{scope}[shift={(12,0)}]
            \fill[color = colorEmphasis, even odd rule] (0,0) circle (0.7) (0.8,0) circle (0.7);
            
            % Vẽ đường viền
            \draw (0,0) circle (0.7);
            \draw (0.8,0) circle (0.7);
            
            % Nhãn
            \node[above left] at (-0.4, 0.4) {$x$};
            \node[above right] at (1.2, 0.4) {$y$};
            \node[below] at (0.4, -0.9) {$x \vartriangle y$};
        \end{scope}

    \end{tikzpicture}
    \caption{Biểu diễn các phép toán tập hợp bằng sơ đồ Ven.}
    \label{fig:li_thuyet_tap_hop:tien_de:tien_de_tach:ven_phep_tinh}
\end{figure}

\exercise Chứng minh rằng, với mọi tập hợp $x$ và $y$, $\bigcap \{x; y\} = x \cap y$.

\solution 

Gọi $z$ là một phần tử trong $\bigcap \{x; y\}$. Khi đó, theo định nghĩa $$\forall a \in \{x; y\}, z \in a.$$
Có $x \in \{x; y\}$ nên $z \in x$, $y \in \{x; y\}$ nên $z \in y$. Do đó, $z \in x \land z \in y$, dẫn đến $z \in x \cup y$. Vậy, $\bigcap \{x; y\} \subseteq x \cap y$.

Ngược lại, giả sử $w \in x \cap y$. Khi này, $w \in x$ và $w \in y$. Nhắc lại rằng $\forall b, b \in \{x; y\} \iff \left[\begin{array}{l}
   b = x \\
   b = y
\end{array}\right.$. Trong cả hai trường hợp đó, $w \in b$. Do vậy, $\forall b \in \{x; y\}, w \in b$. Theo định nghĩa $w \in \bigcap \{x; y\}$. Vậy, $x \cap y \subseteq \bigcap \{x; y\}$.

Kết hợp hai kết quả, chúng ta có $\bigcap \{x; y\} = x \cap y$, theo cách nói khác, điều phải chứng minh.

\exercise Cho $A$, $B$, $C$ là các tập hợp. Chứng minh rằng
\begin{multicols}{2}
   \begin{enumerate}
      \item $\bigcap \{A\} = A$;
      \item $\bigcap \varnothing = \varnothing$
      \item $A \cap B \subseteq A$;
      \item $A \cap B = A \iff A \subseteq B$;
      \item $A \cap B = B \cap A$;
      \item $\bigcap (A; B) = \{A\}$;
      \item $A \subseteq B \implies \bigcap A \supseteq \bigcap B$;
      \item $A \subseteq B \implies A \cap C \subseteq B \cap C$;
      \item $A = B \implies \begin{cases}
               \bigcap A = \bigcap B \\
               A \cap C = B \cap C
            \end{cases}$; 
      \item $(A \cap B) \cap C = A \cap (B \cap C) = \bigcap \{A; B; C\}$;
      \item $\mathfrak{P}(A \cap B) = \mathfrak{P}(A) \cap \mathfrak{P}(B)$.
   \end{enumerate} 
\end{multicols}

\solution

\setcounter{subexercise}{1}
\arabic{subexercise}.Từ định nghĩa,
\begin{align*}
   \bigcap \{A\} &= \left\{z \in \bigcup \{A\} \;\middle|\; \forall a \in \{A\}, z\in a\right\} \\
      &= \left\{z \in A \mid \forall a \in \{A\}, z \in a\right\}.
\end{align*}
Để ý rằng, nếu giả sử $z \in A$, để $a \in \{A\}$ thì $a$ phải bằng $A$, và dẫn đến sự hiển nhiên $z \in a$. Do vậy,
\begin{equation*}
   \forall z, \left(z \in A \implies \forall a \in \{A\}, z \in a\right).
\end{equation*}
Ngoài ra, giả sử $\forall a \in \{A\}, z \in a$. Nhắc lại, $A \in \{A\}$, do đó $z \in A$. Điều này có nghĩa
\begin{equation*}
   \forall z, \left(\forall a \in \{A\}, (z \in a) \implies z \in A\right).
\end{equation*}
Dẫn đến $\forall z, \left(\forall a \in \{A\}, (z \in a) \iff z \in A\right).$ và qua đó $\bigcap \{A\} = \{z \in A \mid z \in A\}$. 

Sử dụng ``$(z) \in A$'' là một vị từ theo hạng thức $z$, từ tiên đề tách, tồn tại tập $y$ để $\forall z, (z \in y \iff z \in y \land z \in A)$. Vậy $y$ có thể là tập nào? Tập $A$ là một ví dụ, do $\forall z, (z \in A \iff z \in A \land z \in A)$. Nhưng vì tính duy nhất của tập xác định từ tiên đề tách, dẫn đến
\begin{equation*}
   \{z \in A \mid z \in A\} = A.
\end{equation*}

Theo tính chất bắc cầu, chúng ta có $\bigcap \{A\} = A$, điều phải chứng minh.

\stepcounter{subexercise}
\arabic{subexercise}
\begin{align*}
   \bigcap \varnothing &= \left\{z \in \bigcup \varnothing \;\middle|\; \forall a \in \varnothing, z\in a\right\} \\ 
   &= \left\{z \in \varnothing \;\middle|\; \forall a \in \varnothing, z\in a\right\}.
\end{align*}
Do đó, với mọi $z$, để $z \in \bigcap \varnothing$ thì trước hết, $z \in \varnothing$. Mặc dù vậy, $z \notin \varnothing$ nên 
\begin{equation*}
   \forall z, z \notin \bigcap \varnothing.
\end{equation*}
Vậy, từ tiên đề tập rỗng, $\bigcap \varnothing = \varnothing$.

\stepcounter{subexercise}
\arabic{subexercise}.Nhắc lại, $A \cap B$ là tập hợp được lấy ra từ tiên đề tách với tập ban đầu là $A$ và vị từ là ``$() \in B$''. Do đó, theo hệ quả \refeq{eq:li_thuyet_tap_hop:tien_de:tien_de_tach:he_qua_tap_con} đã có ở trang \pageref{eq:li_thuyet_tap_hop:tien_de:tien_de_tach:he_qua_tap_con}, $A \cap B \subseteq A$. Điều phải chứng minh.

\stepcounter{subexercise}
\arabic{subexercise}.Giả sử $A \cap B = A$. Xét một phần tử $x \in A$ bất kì. Khi đó, $x \in A \cap B$, và theo định nghĩa của giao, dẫn đến $x \in B$. Do đó, $A \subseteq B$.

Giả sử chiều ngược lại, $A \subseteq B$. Với mọi $e$, nếu $e$ nằm trong $A$ thì $E$ cũng trong $B$, cho chúng ta
$$
\forall e, (e \in A \implies e \in A \land e \in B).
$$
Hiển nhiên có $\forall x, (x \in A \land x \in B \implies x \in A)$. Do vậy, $\forall x, (x \in A \iff x \in A \land x \in B)$. Từ tiên đề tách, có thể thấy $A$ là tập được tách ra từ $A$ với vị từ ``$() \in B$'', nhưng đây cũng là định nghĩa của $A \cap B$. Do đó, $A \cap B = A$.

Kết hợp hai chiều, chúng ta có $A \cap B = A \iff A \subseteq B$, điều phải chứng minh.

\stepcounter{subexercise}
\arabic{subexercise}.Từ tiên đề tách và định nghĩa của giao,
\begin{equation*}
    \forall z, \left(z \in A \cap B \iff z \in A \land z \in B\right).
\end{equation*}
Tương tự, chúng ta cũng có
\begin{equation*}
    \forall z, \left(z \in B \cap A \iff z \in B \land z \in A\right).
\end{equation*}
Theo tính chất bắc cầu của phép đẳng giá, $\forall z, (z \in A \cap B \iff z \in B \cap A)$, và từ tiên đề ngoại diên, $A \cap B = B \cap A$. Điều phải chứng minh.

\stepcounter{subexercise}
\arabic{subexercise}.Nhớ lại, $(A; B) = \{\{A\}; \{A; B\}\}$, dẫn đến $\bigcap (A; B) = \{A\} \cap \{A; B\}$.

Nhận thấy rằng với mọi $w$, để $w$ thuộc $\bigcap (A; B)$ thì $w$ phải thuộc cả $\{A\}$ và $\{A; B\}$. Điều này tương đương với $w = A$. Do vậy, $$\forall w, \left(w \in \bigcap (A; B) \iff w = A\right).$$

Như một kết quả của tiên đề cặp và tiên đề ngoại diên, chỉ tồn tại duy nhất một tập $z$ sao cho $\forall w, (w \in z \iff w = A)$. Do vậy, $\bigcap (A; B) = \{A\}$ và chúng ta có điều phải chứng minh.

\stepcounter{subexercise}
\arabic{subexercise}.Giả sử $A \subseteq B$. Xét $z$ là một phần tử thuộc $\bigcap B$. Khi đó, $z \in b$ với mọi $b \in B$. Do $A \subseteq B$, mọi phần tử $x \in A$ cũng là một phần tử của $B$. Dẫn đến, viết lại các mệnh đề để có
\begin{equation*}
    \begin{cases}
        \forall z \in \bigcap B, \forall b, (b \in B \implies z \in b) \\
        \forall x, (x \in A \implies x \in B)
    \end{cases}.
\end{equation*}
Kết hợp các tính chất của lượng từ rỗng, tính chất phân phối của lượng từ phổ quát với phép hội và tính cụ thể hóa,
\begin{equation*}
    \forall z \in \bigcap B, \forall x, (x \in A \implies z \in x)
\end{equation*}
hay $\forall z, \left(z \in \bigcap B \implies \forall x \in A, z \in x\right)$. 

Ngoài ra, định nghĩa cho chúng ta
$\bigcap A = \{z \in \bigcup A \mid \forall x \in A, z \in x\}$, suy ra \begin{equation*}
    \forall z, \left(z \in \bigcap A \iff z \in \bigcup A \land \forall x \in A, z \in x\right).
\end{equation*}
Nhưng, nhận thấy rằng, với mọi $z$, $\forall x \in A, z \in x \implies \exists x \in A, z \in x$\footnote{Một số người có thể sẽ không chấp nhận cách diễn đạt này với phản biện rằng: ``$A$ có thể là tập rỗng, dẫn đến không tồn tại $x \in A$, và do đó, không thể có $x \in A$ để xét $z \in x$.''. Hãy nhớ lại rằng, theo các khái niệm mà tác giả đã đưa ra, đây chỉ là cách viết tắt. Biểu diễn đầy đủ của mệnh đề đó sẽ là $$
\forall x, (x \in A \implies z \in x) \implies \exists x, (x \in A \implies z \in x)
$$ và điều này thì hoàn toàn đúng.}, dẫn đến $\forall x \in A, z \in x \implies z \in \bigcup A$ là đúng với mọi $z$. Và vì vậy,
\begin{equation}
   \forall z, \begin{cases}
      z \in \bigcap A \iff z \in \bigcup A \land \forall x \in A, z \in x \\
      \forall x \in A, z \in x \implies z \in \bigcup A
   \end{cases}.
   \label{eq:li_thuyet_tap_hop:tien_de:tien_de_tach:7}
\end{equation}

\begin{table}[H]
	\centering
	\caption{Bảng giá trị chân lí của $( P \iff Q \land R ) \land ( R \implies Q ) \implies ( P \iff R )$.}
   \label{tab:li_thuyet_tap_hop:tien_de:tien_de_tach:7}
	\begin{tblr}{
		colspec = {|c|c|c|ccccccccccccc|}
	}
		\hline
		$P$ & $Q$ & $R$ & $(P$ & $\iff$ & $Q$ & $\land$ & $R)$ & $\land$ & $(R$ & $\implies$ & $Q)$ & $\implies$ & $(P$ & $\iff$ & $R)$ \\
		\hline[1.5pt]
		Đ & Đ & Đ & Đ & Đ & Đ & Đ & Đ & Đ & Đ & Đ & Đ & \emphcolor{Đ} & Đ & Đ & Đ \\
		Đ & Đ & S & Đ & S & Đ & S & S & S & S & Đ & Đ & \emphcolor{Đ} & Đ & S & S \\
		Đ & S & Đ & Đ & S & S & S & Đ & S & Đ & S & S & \emphcolor{Đ} & Đ & Đ & Đ \\
		Đ & S & S & Đ & S & S & S & S & S & S & Đ & S & \emphcolor{Đ} & Đ & S & S \\
		S & Đ & Đ & S & S & Đ & Đ & Đ & S & Đ & Đ & Đ & \emphcolor{Đ} & S & S & Đ \\
		S & Đ & S & S & Đ & Đ & S & S & Đ & S & Đ & Đ & \emphcolor{Đ} & S & Đ & S \\
		S & S & Đ & S & Đ & S & S & Đ & S & Đ & S & S & \emphcolor{Đ} & S & S & Đ \\
		S & S & S & S & Đ & S & S & S & Đ & S & Đ & S & \emphcolor{Đ} & S & Đ & S \\
		\hline
	\end{tblr}
\end{table}

Bảng giá trị chân lí \ref{tab:li_thuyet_tap_hop:tien_de:tien_de_tach:7} cho phép chúng ta suy ra từ \refeq{eq:li_thuyet_tap_hop:tien_de:tien_de_tach:7}:
\begin{equation*}
   \forall z, \left(z \in \bigcap A \iff \forall x \in A, z\in x\right).
\end{equation*}

Sử dụng kết quả này để thay thế tương đương, có $\forall z, (z \in \bigcap B \implies z \in \bigcap A)$. Qua đó, $\bigcap B \subseteq \bigcap A$. Từ đây, chúng ta có thể dễ dàng nhìn ra được điều phải chứng minh.

\stepcounter{subexercise}
\arabic{subexercise}. Giả sử $A \subseteq B$. Gọi $x \in A \cap C$ bất kì. Chúng ta có đồng thời $x \in A$ và $x \in C$. Từ giả thiết ban đầu, $x \in A$ có được $x \in B$. Và khi chúng ta có cả $x \in B$ và $x \in C$, $x \in B \cap C$.

Tóm lại, $x \in A \cap C \implies x \in B \cap C$ với mọi $x$, cho nên $A \cap C \subseteq B \cap C$. Qua đó, chúng ta có điều phải chứng minh.

\stepcounter{subexercise}
\arabic{subexercise}. Chúng ta đã chứng minh:
\begin{equation*}
   A \subseteq B \implies
   \begin{cases}
      \bigcap A \supseteq \bigcap B \\
      A \cap C \subseteq B \cap C
   \end{cases}.
\end{equation*}
Và từ đó, chúng ta cũng sẽ có
\begin{equation*}
   A \supseteq B \implies
   \begin{cases}
      \bigcap A \subseteq \bigcap B \\
      A \cap C \supseteq B \cap C
   \end{cases}.
\end{equation*}

Có $x = y \iff (x \subseteq y) \land (x \supseteq y)$ đúng với mọi tập $x$ và $y$. Cho nên, chúng ta dễ dàng thấy được
\begin{equation*}
   A = B \implies \begin{cases}
               \bigcap A = \bigcap B \\
               A \cap C = B \cap C
            \end{cases}.
\end{equation*}
Điểu phải chứng minh.
   
\stepcounter{subexercise}
\arabic{subexercise}. Trước hết, chúng ta sẽ chứng minh $(A \cap B) \cap C = A \cap (B \cap C)$. Thật vậy, với mọi $x$, sử dụng liên tục định nghĩa giao hai tập hợp, chúng ta có
\begin{align*}
   x \in (A \cap B) \cap C &\iff \begin{cases}
      x \in A \cap B \\
      x \in C
   \end{cases}; \\
   x \in (A \cap B) \cap C &\iff \begin{cases}
      x \in A \\
      x \in B \\
      x \in C
   \end{cases}; \\
   x \in (A \cap B) \cap C &\iff \begin{cases}
      x \in A \\
      x \in B \cap C
   \end{cases}; \\
   x \in (A \cap B) \cap C &\iff x \in A \cap (B \cap C).
\end{align*}

Tiếp theo, chúng ta sẽ cần phải khẳng định rằng $(A \cap B) \cap C = \bigcap \{A; B; C\}$. Lấy tùy ý $z \in \bigcap\{A; B; C\}$, điều này tương đương với $z \in \bigcup \{A; B; C\}$ và $\forall y \in \{A; B; C\}, z \in y$. Mà, dễ nhận thấy rằng $\forall y \in \{A; B; C\}, z \in y \implies z \in \bigcup \{A; B; C\}$, cho nên
\begin{equation*}
   z \in \bigcap\{A; B; C\} \iff \forall y \in \{A; B; C\}, z \in y.
\end{equation*}

$\{A; B; C\}$ là tập duy nhất thỏa mãn với mọi $y$, $\left[\begin{array}{l}
   y = A \\
   y = B \\
   y = C   
\end{array}\right.$. Do vậy, với mọi $z$,
\begin{align*}
   \forall y \in \{A; B; C\}, z \in y &\iff \forall y, \left(\left[\begin{array}{l}
   y = A \\
   y = B \\
   y = C   
\end{array} \implies z\in y\right.\right); \\
\forall y \in \{A; B; C\}, z \in y &\iff \forall y, \begin{cases}
   y = A \implies z \in y \\
   y = B \implies z \in y \\
   y = C \implies z \in y 
\end{cases} \equationexplanation{quy tắc chia trường hợp}; \\
\forall y \in \{A; B; C\}, z \in y &\iff \forall y, \begin{cases}
   z \in A \\
   z \in B \\
   z \in C
\end{cases}.
\end{align*}
Bỏ lượng từ rỗng $\forall y$ đi để được $\forall y \in \{A; B; C\}, z \in y \iff \begin{cases}
   z \in A \\
   z \in B \\
   z \in C
\end{cases}$. Đã có $z \in (A \cap B) \cap C \iff \begin{cases}
   z \in A \\
   z \in B \\
   z \in C
\end{cases}$, cho nên bắc cầu để có $$\forall z, z \in \bigcap\{A; B; C\} \iff z \in (A \cap B) \cap C.$$
Do vậy $\bigcap\{A; B; C\} = (A \cap B) \cap C$. Từ đây, theo tính chất bắc cầu của \coloringCrimson{=}, dễ thấy điều phải chứng minh.

\stepcounter{subexercise}
\arabic{subexercise}. Vọi mọi $x$,
\begin{align*}
   x \in \mathfrak{P}(A \cap B) &\iff x \subseteq A \cap B \equationexplanation{định nghĩa tập lũy thừa}; \\
   x \in \mathfrak{P}(A \cap B) &\iff \forall y, (y \in x \implies y \in A \cap B) \equationexplanation{định nghĩa tập con}; \\
   x \in \mathfrak{P}(A \cap B) &\iff \forall y, (y \in x \implies y \in A \land y \in B) \equationexplanation{định nghĩa giao hai tập hợp}; \\
   x \in \mathfrak{P}(A \cap B) &\iff \forall y, \begin{cases}
      y \in x \implies y \in A \\
      y \in x \implies y \in B
   \end{cases} \equationexplanation{$(P \implies (Q \land R)) \iff \begin{cases}
      P \implies Q \\
      P \implies R
   \end{cases}$}; \\
   x \in \mathfrak{P}(A \cap B) &\iff \begin{cases}
      \forall y, y \in x \implies y \in A \\
      \forall y, y \in x \implies y \in B
   \end{cases} \equationexplanation{\begin{tabular}{l}
      tính chất phân phối của lượng từ \\
      phổ quát với phép hội
   \end{tabular}}; \\
   x \in \mathfrak{P}(A \cap B) &\iff x \subseteq A \land x \subseteq B \equationexplanation{định nghĩa tập con}; \\
   x \in \mathfrak{P}(A \cap B) &\iff x \in \mathfrak{P}(A) \land x \in \mathfrak{P}(B) \equationexplanation{định nghĩa tập lũy thừa}; \\
   x \in \mathfrak{P}(A \cap B) &\iff x \in \mathfrak{P}(A) \cap \mathfrak{P}(B) \equationexplanation{định nghĩa giao hai tập hợp}.
\end{align*} 
Và từ đó, $\mathfrak{P}(A \cap B) = \mathfrak{P}(A) \cap \mathfrak{P}(B)$. Điều phải chứng minh.

\exercise Cho $A, B, C, D, E$ là các tập hợp. Chứng minh rằng
\begin{enumerate}
   \item $A \cap (B \cup C) = (A \cap B) \cup (A \cap C)$;
   \item $A \cup (B \cap C) = (A \cup B) \cap (A \cup C)$;
   \item $A \cup B = A \cap C \iff \begin{cases}
      B \subseteq A \\
      A \subseteq C
   \end{cases}$;
   \item $\begin{cases}
      A \cap B = A \cap C \\
      A \cup B = A \cup C
   \end{cases} \iff B = C$;
   \item $A \cap (B \cup (A \cap C)) = A \cap (B \cup C)$;
   \item $A \cup (B \cap (A \cup C)) = A \cup (B \cap C)$;
   \item $A \cap B = C \cap D \implies (A \cup (B \cap C)) \cap (A \cup (B \cap D)) = A$;
   \item $\begin{cases}
      A \cap B = C\cap D\\
      C \cup D = E \\
      C \subseteq A \\
      D \subseteq B \\
      A \subseteq E \\
      B \subseteq E
   \end{cases} \implies \begin{cases}
      C = A\\
      D = B
   \end{cases}$.
\end{enumerate}
\solution

\setcounter{subexercise}{1}
\arabic{subexercise}. Từ những kết quả đã có, với mọi $x$, thực hiện biến đổi như sau
\begin{align*}
   x \in A \cap (B \cup C) &\iff x \in A \land x \in B \cup C; \\
   x \in A \cap (B \cup C) &\iff x \in A \land (x \in B \lor x \in C); \\
   x \in A \cap (B \cup C) &\iff (x \in A \land x \in B) \lor (x \in A \land x \in C) \equationexplanation{tính chất phân phối}; \\
   x \in A \cap (B \cup C) &\iff x \in A \cap B \lor x \in A \cap C; \\
   x \in A \cap (B \cup C) &\iff x \in (A \cap B) \cup (A \cap C).
\end{align*}
Từ đó, dễ có điều phải chứng minh.

\stepcounter{subexercise}
\arabic{subexercise}. Tương tự, với mọi $x$,
\begin{align*}
   x \in A \cup (B \cap C) &\iff x \in A \lor x \in B \cap C; \\
   x \in A \cup (B \cap C) &\iff x \in A \lor (x \in B \land x \in C); \\
   x \in A \cup (B \cap C) &\iff (x \in A \lor x \in B) \land (x \in A \lor x \in C); \\
   x \in A \cup (B \cap C) &\iff x \in A \cup B \land x \in A \cup C; \\
   x \in A \cup (B \cap C) &\iff x \in (A \cup B) \cap (A \cup C)
\end{align*}
và điều phải chứng minh là hiển nhiên.

\stepcounter{subexercise}
\arabic{subexercise}. Đặt $S = A \cup B = A \cap C$. Gọi $x$ là một phần tử thuộc $S$. Theo cách đặt, $x$ cũng thuộc $A \cap C$, dẫn đến $x \in A$. Qua đó, chúng ta được
$$\forall x, x \in S \implies x \in A.$$

Giả sử $x \in A$. Dễ thấy $x \in A \cup B$ cũng đúng, và từ đó, $x \in S$. Do vậy, chúng ta cũng có $$\forall x, x \in A \implies x \in S.$$

Suy ra, $\forall x, x \in A \iff x \in S$. Vậy, $S = A$, cho chúng ta $A \cup B = A \cap C \iff \begin {cases}A = A \cup B \\ A = A \cap C\end{cases}$. Dùng những bài tập trước như hệ quả,
\begin{equation*}
   \begin{cases}
      A = A \cup B \iff B \subseteq A \\
      A = A \cap C \iff A \subseteq C
   \end{cases}.
\end{equation*}
Và vậy, $A \cup B = A \cap C \iff \begin{cases}
   B \subseteq A \\
   A \subseteq C
\end{cases}$. Điều phải chứng minh.

\stepcounter{subexercise}
\arabic{subexercise}. Từ $B = C$ dễ dàng suy ra $\begin{cases}
   A \cap B = A \cap C \\
   A \cup B = A \cup C
\end{cases}$. Chứng minh chiều ngược lại, giả sử $A \cap B = A \cap C$ và $A \cup B = A \cup C$. Khi này,
\begin{align*}
   B  &= (A \cap B) \cup B \\
      &= (A \cap C) \cup B \equationexplanation{giả thiết có $A \cap B = A \cap C$}; \\
      &= (A \cup B) \cap (C \cup B); \\
      &= (C \cup A) \cap (C \cup B) \equationexplanation{$A \cup B = A \cup C = C \cup A$}; \\
      &= C \cup (A \cap B); \\
      &= C \cup (A \cap C) \equationexplanation{$A \cap B = A\cap C$}; \\
      &= C.
\end{align*}
Qua đó, dễ thấy điều phải chứng minh.

\stepcounter{subexercise}
\arabic{subexercise}. Có $A \subseteq B \cup A$ nên $A \cap (B \cup A) = A$. Thực hiện biến đổi để có
\begin{align*}
   A \cap (B \cup (A \cap C)) &= A \cap ((B \cup A) \cap (B \cup C)) \\
   &= (A \cap (B \cup A)) \cap (B \cup C) \\
   &= A \cap (B \cup C).
\end{align*}
Điều phải chứng minh.

\stepcounter{subexercise}
\arabic{subexercise}. Làm tương tự, với để ý rằng $A \supseteq B \cap A$ nên $A \cup (B \cap A) = A$,
\begin{align*}
   A \cup (B \cap (A \cup C)) &= A \cup ((B \cap A) \cup (B \cap C)) \\
   &= (A \cup (B \cap A)) \cup (B \cap C) \\
   &= A \cup (B \cap C).
\end{align*}
Điều phải chứng minh.

\stepcounter{subexercise}
\arabic{subexercise}. Giả sử $A \cap B = C \cap D$,
\begin{align*}
   (A \cup (B \cap C)) \cap (A \cup (B \cap D)) &= A \cup ((B \cap C) \cap (B \cap D)) \\
   &= A \cup (((B \cap C) \cap B) \cap D) \\
   &= A \cup ((B \cap C) \cap D) \equationexplanation{$B \cap C \subseteq B$ nên $(B \cap C) \cap B = B \cap C$} \\
   &= A \cup (B \cap (C \cap D)) \\
   &= A \cup (B \cap (A \cap B)) \equationexplanation{giả thiết} \\
   &= A \cup (A \cap B) \equationexplanation{$A \cap B \subseteq B$ nên $B \cap (A \cap B) = A \cap B$} \\
   &= A.
\end{align*}
Điều phải chứng minh.

\stepcounter{subexercise}
\arabic{subexercise}. Lấy tùy ý $y \in A$. Do $A \subseteq E$ nên $y \in E$, hay $y \in C \cup D$. Suy ra, $y \in C$ hoặc $y \in D$.

Giả sử $y \notin C$, dẫn đến $y \in D$. Kết hợp $D \subseteq B$ để có $y \in B$. $y$ đồng thời thuộc $A$ và $B$ cho khẳng định $y \in A \cap B$. Đề bài cho $A \cap B = C\cap D$ nên $y \in C \cap D$. Vậy $y \in C$, nhưng điều này mâu thuẫn với giả sử ban đầu.

Sử dụng phương pháp chứng minh bằng phản chứng, cần phải có $y\in A \implies y \in C$ với mọi $y$, và đương nhiên, khi đó, $A \subseteq C$. Đã có $C \subseteq A$ nên $C = A$.

Với lập luận tương đương, chúng ta cũng có thể suy ra $D = B$. Kết hợp lại để có được điều phải chứng minh.

\exercise Cho $A, B, C$ là các tập hợp. Chứng minh rằng
\begin{multicols}{2}
   \begin{enumerate}
      \item $A \setminus B \subseteq A$;
      \item $A \setminus B = A \iff A \cap B = \varnothing$;
      \item $A \setminus B = \varnothing \iff A \subseteq B$;
      \item $A \setminus B = A \setminus (A \cap B)$;
      \item $A = (A \setminus B) \cup (A \cap B)$;
      \item $A \subseteq B \iff \begin{cases}
         A \setminus C \subseteq B \setminus C \\
         C \setminus A \supseteq C \setminus B
      \end{cases}$;
      \item $A = B \iff \begin{cases}
         A \setminus C = B \setminus C \\
         C \setminus A = C \setminus B
      \end{cases}$;
      \item $A \setminus (B\cap C) = (A \setminus B) \cup (A \setminus C)$;
      \item $A \setminus (B \cup C) = (A \setminus B) \cap (A \setminus C) = (A \setminus B) \setminus C$;
      \item $A \setminus (B \setminus C) = (A \setminus B) \cup (A \cap C)$;
      \item $(A \setminus B) \setminus (A \setminus C) = (A \setminus B) \cap C = (A \cap C) \setminus B$;
      \item $A \subseteq B \iff \begin{cases}
         A \cap C \subseteq B \cap C \\
         A \setminus C \subseteq B \setminus C
      \end{cases}$.
   \end{enumerate}
\end{multicols}

\solution

\setcounter{subexercise}{1}
\arabic{subexercise}. Có $A \setminus B = \{x \in A \mid x \notin B\}$. Theo hệ quả của tiên đề tách \refeq{eq:li_thuyet_tap_hop:tien_de:tien_de_tach:he_qua_tap_con}, $A \setminus B \subseteq A$. Điều phải chứng minh.

\stepcounter{subexercise}
\arabic{subexercise}. Đầu tiên, giả sử $A \setminus B = A$. Xét một phần tử $x$.
\begin{itemize}
   \item Nếu $x \notin A$, hiển nhiên $x \notin A \cap B$;
   \item Nếu $x \in A$, do $A = A \setminus B$ nên $x \notin B$. Qua đó, $x \notin A \cap B$.
\end{itemize}
Do $x$ chỉ có thể thuộc hoặc không thuộc $A$, cho nên chúng ta có thể suy ra rằng $\forall x, x \notin A \cap B$. Định nghĩa của tập rỗng cũng có biểu thức tương tự. Chỉ có một tập rỗng nên $A \cap B = \varnothing$.

Chứng minh chiều ngược lại, cho $A \cap B = \varnothing$. Gọi $y$ là một phần tử bất kì thuộc $A$. Do $y \notin A \cap B$ (tính chất tập rỗng) nên không thể xảy ra đồng thời $y \in A$ và $y \in B$. Mà $y \in A$ theo giả thiết, dẫn đến $y \notin B$. Qua đó, $y \in A \setminus B$ với mọi $y$ trong $A$, suy ra $A \subseteq A\setminus B$.

Từ phần trước của bài tập này, chúng ta có $A\setminus B \subseteq A$, nên phải có $A = A \setminus B$. 

Kết hợp hai chiều, chúng ta có điều phải chứng minh.

\stepcounter{subexercise}
\arabic{subexercise}. Giả sử $A \setminus B = \varnothing$. Lấy $x \in A$ bất kì. Nếu $x \notin B$, sẽ dẫn đến $x \in A \setminus B$. Nhưng, theo tiên đề tập rỗng, do $A \setminus B = \varnothing$ nên $\forall x, x \notin A \setminus B$, mâu thuẫn. Do vậy, $x \in B$. Chúng ta đã chứng minh, $$\forall x \in A, x \in B$$ hay $A \subseteq B$ trong trường hợp này.

Trường hợp còn lại, giả sử $A \subseteq B$. Xét một phần tử $y$ bất kì.
\begin{itemize}
   \item Nếu $y \notin A$, hiển nhiên $y \notin A \setminus B$;
   \item Nếu $y \in A$, dễ thấy $y$ cũng phải thuộc $B$. Qua đó, $y \notin A \setminus B$.
\end{itemize}
Vậy, với mọi $y$, $y$ luôn không ở trong $A \setminus B$. Hay, $A \setminus B = \varnothing$.

Kết hợp hai chiều, điều phải chứng minh.

\stepcounter{subexercise}
\arabic{subexercise}. Theo định nghĩa của phép hiệu và phép giao:
\begin{align*}
x \in A \setminus (A \cap B) &\iff (x \in A) \land \overline{x \in A \cap B}; \\
x \in A \setminus (A \cap B) &\iff (x \in A) \land \overline{x \in A \land x \in B}; \\
x \in A \setminus (A \cap B) &\iff (x \in A) \land (x \notin A \lor x \notin B) \equationexplanation{định luật Đờ Moóc-gơn}; \\
x \in A \setminus (A \cap B) &\iff ((x \in A) \land (x \notin A)) \lor ((x \in A) \land (x \notin B)) \equationexplanation{tính chất phân phối}.
\end{align*}

Do một phần tử không thể vừa thuộc vừa không thuộc một tập hợp, cho nên
\begin{align*}
   x \in A \setminus (A \cap B) &\iff x \in A \land x \notin B; \\
   x \in A \setminus (A \cap B) &\iff x \in A \setminus B.
\end{align*}

Vậy, $A \setminus (A \cap B) = A \setminus B$. Điều phải chứng minh.

\stepcounter{subexercise}
\arabic{subexercise}. Lấy một phần tử $x$ bất kì thuộc $A$. Xét hai trường hợp xảy ra đối với tập hợp $B$:
\begin{itemize}
    \item Nếu $x \in B$, vì $x \in A$ và $x \in B$, ta có $x \in A \cap B$.\\
    Suy ra $x \in (A \setminus B) \cup (A \cap B)$.
    
    \item Nếu $x \notin B$, vì $x \in A$ và $x \notin B$, ta có $x \in A \setminus B$.\\
    Suy ra $x \in (A \setminus B) \cup (A \cap B)$.
\end{itemize}
Trong cả hai trường hợp, đều có $x \in (A \setminus B) \cup (A \cap B)$.\\
Do đó,
\[ A \subseteq (A \setminus B) \cup (A \cap B). \]

Lấy một phần tử $y \in (A \setminus B) \cup (A \cap B)$. Theo định nghĩa phép hợp,
\begin{itemize}
   \item $x \in A \setminus B$. Theo định nghĩa phép hiệu, $x \in A$ (và $x \notin B$);  
   \item Hoặc $x \in A \cap B$. Theo định nghĩa phép giao, $x \in A$ (và $x \in B$).
\end{itemize}
Trong cả hai trường hợp, chúng ta đều có $x \in A$.\\
Vì vậy,
\[ (A \setminus B) \cup (A \cap B) \subseteq A. \]

Vậy hai tập hợp là bằng nhau, điều phải chứng minh.

\stepcounter{subexercise}
\arabic{subexercise}. Trước hết, giả sử $A \subseteq B$. Lấy một phần tử $x$ bất kì thuộc vào tập hợp $A \setminus C$. Theo định nghĩa của phép hiệu tập hợp,
$ x \in A \land x \notin C. $
Vì $A \subseteq B$ và $x \in A$, nên
$ x \in B. $
Từ đó, có đồng thời:
\[ x \in B \land x \notin C. \]
Theo định nghĩa của phép hiệu, điều này có nghĩa là $ x \in B \setminus C $
Vì mọi phần tử $x \in A \setminus C$ đều thuộc $B \setminus C$, kết luận được
\[ A \setminus C \subseteq B \setminus C. \]

Lấy một phần tử $y$ bất kì thuộc vào tập hợp $C \setminus B$. Chúng ta có $ y \in C \land y \notin B $. Cần chứng minh rằng $y \notin A$. Giả sử ngược lại rằng $y \in A$. Vì $A \subseteq B$, nếu $y \in A$ thì ta bắt buộc phải có $y \in B$. Điều này mâu thuẫn trực tiếp với điều kiện ban đầu là $y \notin B$.
Do đó, giả thiết phản chứng là sai. Dẫn đến $ y \notin A $. Như vậy, chúng ta có đồng thời
\[ y \in C \land y \notin A. \]
Theo định nghĩa của phép hiệu, điều này có nghĩa là $ y \in C \setminus A $. Vậy
\[ C \setminus B \subseteq C \setminus A \quad \text{hay} \quad C \setminus A \supseteq C \setminus B. \]

Kết hợp hai kết quả, chúng ta có 
\begin{equation}
   A \subseteq B \implies \begin{cases}
      A \setminus C \subseteq B \setminus C \\
      C \setminus A \supseteq C \setminus B
   \end{cases}.
   \label{eq:li_thuyet_tap_hop:tien_de:tien_de_tach:so_sanh_hieu.1}
\end{equation}

Sau đó, chúng ta giả sử chiều ngược lại rằng $\begin{cases}
   A \setminus C \subseteq B \setminus C \\
   C \setminus A \supseteq C \setminus B
\end{cases}$. Sử dụng chứng minh bằng phương pháp phản chứng, giả sử $A \not \subseteq B$, tương đương với $\overline{\forall z, (z \in A \implies z \in B)}$. Theo lập luận lô-gích thông thường, $\exists z, (z \in A \land z \notin B)$. Cụ thể hóa phần tử đó và gọi nó là $w$. Xét hai trường hợp:
\begin{itemize}
   \item Trường hợp \coloringCrimson{$w \in C$}: Khi này, chúng ta có $w \in C \land w \in A$ và $w \in C \land w \notin B$. Dẫn đến
   \begin{equation*}
      \begin{cases}
         w \notin C \setminus A\\
         w \in C \setminus B
      \end{cases}.
   \end{equation*} Chúng ta thấy tồn tại phần tử $w$ trong $C \setminus B$ không thuộc $C \setminus A$, nhưng giả thiết $C \setminus A \supseteq C \setminus B$ lại mâu thuẫn với điều này.
   \item Trường hợp \coloringDodgerblue{$w \notin C$}: Chúng ta cũng dễ dàng suy ra được $w \in A \setminus C$ và $w \notin B \setminus C$, nhưng $A \setminus C \subseteq B \setminus C$ cũng tạo ra mâu thuẫn.
\end{itemize}
Do mọi trường hợp đều ra mâu thuẫn, nên điều giả sử ban đầu là sai. Do đó, $A \subseteq B$. Vậy,
\begin{equation}
   \begin{cases}
      A \setminus C \subseteq B \setminus C \\
      C \setminus A \supseteq C \setminus B
   \end{cases} \implies A \subseteq B.
   \label{eq:li_thuyet_tap_hop:tien_de:tien_de_tach:so_sanh_hieu.2}
\end{equation}

Từ \refeq{eq:li_thuyet_tap_hop:tien_de:tien_de_tach:so_sanh_hieu.1} và \refeq{eq:li_thuyet_tap_hop:tien_de:tien_de_tach:so_sanh_hieu.2}, chúng ta có điều phải chứng minh.

\stepcounter{subexercise}
\arabic{subexercise}. Phần trước cho 
\begin{equation*}
   A \subseteq B \iff \begin{cases}
      A \setminus C \subseteq B \setminus C \\
      C \setminus A \supseteq C \setminus B
   \end{cases}.
\end{equation*} Sử dụng chính kết quả này, chúng ta cũng có
\begin{equation*}
   B \subseteq A \iff \begin{cases}
      B \setminus C \subseteq A \setminus C \\
      C \setminus B \subseteq C \setminus A
   \end{cases}.
\end{equation*}
Qua đó,
\begin{equation*}
   A \subseteq B \land B \subseteq A \iff \begin{cases}
      A \setminus C \subseteq B \setminus C \\
      C \setminus A \supseteq C \setminus B \\
      B \setminus C \subseteq A \setminus C \\
      C \setminus B \subseteq C \setminus A
   \end{cases}
\end{equation*}
và dễ dàng suy ra được
\begin{equation*}
   A = B \iff \begin{cases}
      A \setminus C = B \setminus C \\
      C \setminus A = C \setminus B
   \end{cases}.
\end{equation*}
Điều phải chứng minh.

\stepcounter{subexercise}
\arabic{subexercise}. Lấy $x \in A \setminus (B \cap C)$. Khi này,
\begin{align*}
   x \in A \setminus (B \cap C) &\iff x \in A \land x \notin B \cap C; \\
   x \in A \setminus (B \cap C) &\iff x \in A \land \overline{x \in B \cap C}; \\
   x \in A \setminus (B \cap C) &\iff x \in A \land \overline{x \in B \land x \in C}; \\
   x \in A \setminus (B \cap C) &\iff x \in A \land (x \notin B \lor x \notin C) \equationexplanation{định luật Đờ Moóc-gan}; \\
   x \in A \setminus (B \cap C) &\iff (x \in A \land x \notin B) \lor (x \in A \land x\notin C) \equationexplanation{tính chất phân phối}; \\
   x \in A \setminus (B \cap C) &\iff x \in A \setminus B \lor x \in A \setminus C; \\
   x \in A \setminus (B \cap C) &\iff x \in (A \setminus B) \cup (A \setminus C).
\end{align*}
Từ đó, dễ thấy điều phải chứng minh.

\stepcounter{subexercise}
\arabic{subexercise}. Đầu tiên, cần phải chứng minh $A \setminus (B \cup C) = (A \setminus B) \cap (A \setminus C)$. Với lập luận tương tự như phần trước, với mọi $x$,
\begin{align*}
   x \in A \setminus (B \cup C) &\iff x \in A \land x \notin B \cup C; \\
   x \in A \setminus (B \cup C) &\iff x \in A \land \overline{x \in B \cup C}; \\
   x \in A \setminus (B \cup C) &\iff x \in A \land \overline{x \in B \lor x \in C}; \\
   x \in A \setminus (B \cup C) &\iff x \in A \land (x \notin B \land x \notin C) \equationexplanation{định luật Đờ Moóc-gan}; \\
   x \in A \setminus (B \cup C) &\iff (x \in A \land x \notin B) \land (x \in A \land x\notin C); \\
   x \in A \setminus (B \cup C) &\iff x \in A \setminus B \land x \in A \setminus C; \\
   x \in A \setminus (B \cup C) &\iff x \in (A \setminus B) \cap (A \setminus C)
\end{align*}
sẽ dẫn đến $A \setminus (B \cup C) = (A \setminus B) \cap (A \setminus C)$ như chúng ta cần.

Chúng ta cũng có thể lập luận theo một hướng khác:
\begin{align*}
   x \in A \setminus (B \cup C) &\iff x \in A \land (x \notin B \land x \notin C); \\
   x \in A \setminus (B \cup C) &\iff (x \in A \land x \notin B) \land x \notin C; \\
   x \in A \setminus (B \cup C) &\iff x \in A \setminus B \land x \notin C; \\
   x \in A \setminus (B \cup C) &\iff x \in (A \setminus B) \setminus C.
\end{align*}
Và từ đó, chúng ta cũng có $A \setminus (B \cup C) = (A \setminus B) \setminus C$. Bắc cầu hai kết quả để có điều phải chứng minh.

\stepcounter{subexercise}
\arabic{subexercise}. Với mọi $x$,
\begin{align*}
   x \in A \setminus (B \setminus C) &\iff x \in A \land x \notin B \setminus C; \\
   x \in A \setminus (B \setminus C) &\iff x \in A \land \overline{x \in B \land x \notin C}; \\
   x \in A \setminus (B \setminus C) &\iff x \in A \land \left(x \notin B \lor x \in C\right); \\
   x \in A \setminus (B \setminus C) &\iff (x \in A \land x \notin B) \lor (x \in A \land x \in C); \\
   x \in A \setminus (B \setminus C) &\iff x \in A \setminus B \lor x \in A \cap C; \\
   x \in A \setminus (B \setminus C) &\iff x \in (A \setminus B) \cup (A \cap C).
\end{align*}
Do đó, $A \setminus (B \setminus C) = (A \setminus B) \cup (A \cap C)$. Điều phải chứng minh.

\stepcounter{subexercise}
\arabic{subexercise}. Với mọi $x$,
\begin{align*}
   x \in (A \setminus B) \setminus (A \setminus C) &\iff x \in A \setminus B \land x \notin A \setminus C; \\
   x \in (A \setminus B) \setminus (A \setminus C) &\iff (x \in A \land x \notin B) \land \overline{x \in A \land x \notin C}; \\
   x \in (A \setminus B) \setminus (A \setminus C) &\iff (x \in A \land x \notin B) \land \overline{x \in A \land x \notin C}; \\
   x \in (A \setminus B) \setminus (A \setminus C) &\iff x \in A \land x \notin B \land (x \notin A \lor x \in C); \\
   x \in (A \setminus B) \setminus (A \setminus C) &\iff (x \in A \land x \notin B \land x \notin A) \lor (x \in A \land x \notin B \land x \in C).
\end{align*}
Nhận thấy rằng $x \in A \land x \notin A$ là luôn sai để không có mâu thuẫn, cho nên 
\begin{align*}
   x \in (A \setminus B) \setminus (A \setminus C) &\iff x \in A \land x \notin B \land x \in C; \\
   x \in (A \setminus B) \setminus (A \setminus C) &\iff \begin{cases}
      (x \in A \land x \notin B) \land x \in C \\
      (x \in A \land x \in C) \land x \notin B
   \end{cases}; \\
   x \in (A \setminus B) \setminus (A \setminus C) &\iff \begin{cases}
      x \in A \setminus B \land x \in C \\
      x \in A \cap C \land x \notin B
   \end{cases}; \\
   x \in (A \setminus B) \setminus (A \setminus C) &\iff \begin{cases}
      x \in (A \setminus B) \cap C \\
      x \in (A \cap C) \setminus B
   \end{cases}.
\end{align*}
Từ đó, điều phải chứng minh là khá hiển nhiên.

\stepcounter{subexercise}
\arabic{subexercise}. Chứng minh chiều thuận từ $A \subseteq B$ sang $\begin{cases}
   A \cap C \subseteq B \cap C \\
   A \setminus C \subseteq B \setminus C
\end{cases}$ đã được tác giả trình bày ở những bài tập trước. Chứng minh chiều ngược lại, chúng ta giả thiết đã có $\begin{cases}
   A \cap C \subseteq B \cap C \\
   A \setminus C \subseteq B \setminus C
\end{cases}$. Sử dụng chứng minh bằng phương pháp phản chứng, giả sử $A \not \subseteq B$, tương đương với $\exists z, (z \in A \land z \notin B)$. Gọi phần tử thỏa mãn là $w$. Xét hai trường hợp:
\begin{itemize}
   \item Trường hợp \coloringCrimson{$w \in C$}: Khi này, chúng ta có $w \in A \land w \in C$ và $w \notin B \land w \in C$. Dẫn đến
   \begin{equation*}
      \begin{cases}
         w \in A \cap C\\
         w \notin B \cap C
      \end{cases},
   \end{equation*} nhưng giả thiết $A \cap C \supseteq B \cap C$ lại mâu thuẫn với điều này.
   \item Trường hợp \coloringDodgerblue{$w \notin C$}: Chúng ta cũng dễ dàng suy ra được $w \in A \setminus C$ và $w \notin B \setminus C$, nhưng $A \setminus C \subseteq B \setminus C$ cũng tạo ra mâu thuẫn.
\end{itemize}
Do mọi trường hợp đều ra mâu thuẫn, nên điều giả sử ban đầu là sai. Do đó, cần phải có $A \subseteq B$.

Kết hợp hai chiều, chúng ta có điều phải chứng minh.

\exercise Cho $A, B, C$ là các tập hợp. Chứng minh rằng
\begin{multicols}{2}
   \begin{enumerate}
      \item $A \vartriangle B = (A \cup B) \setminus (A \cap B)$;
      \item $A \vartriangle B = B \vartriangle A$;
      \item $A \vartriangle B = A \iff B = \varnothing$;
      \item $A \vartriangle B = \varnothing \iff B = A$;
      \item $A \vartriangle B = C \implies \begin{cases}
         B \vartriangle C = A \\
         C \vartriangle A = B
      \end{cases}$;
      \item $A = B \implies A \vartriangle C = B \vartriangle C$;
      \item $(A \vartriangle B) \vartriangle C = A \vartriangle (B \vartriangle C)$.
   \end{enumerate}
\end{multicols}

\solution

\setcounter{subexercise}{1}
\arabic{subexercise}. Gọi $x \in A \vartriangle B$. Theo định nghĩa, $x \in (A \setminus B) \cup (B \setminus A)$, hay $x \in A \setminus B$ hoặc $x \in B \setminus A$. Qua đó, $x$ cần phải thuộc $A$ hoặc $B$. Do đó $x \in A\cup B$.

Nếu $x \in A\cap B$, điều này có nghĩa là đồng thời $x \in A$ và $x \in B$, dẫn đến $x \in A \setminus B$ và $x \in B \setminus A$ đều không thể xảy ra, mâu thuẫn. Do vậy, $x \notin A \cap B$. Và từ đó, $x \in (A \cup B) \setminus (A \cap B)$ là điều hiển nhiên. Vậy,
\begin{equation*}
   A \vartriangle B \subseteq (A \cup B) \setminus (A \cap B).
\end{equation*}

Ngược lại, với mọi $y$, nếu $y \in (A \cup B) \setminus (A \cap B)$. Chúng ta có $y \in A \cup B$.
\begin{itemize}
   \item Nếu $y \in A$, thấy được rằng $y \notin A \cap B$ nên cần phải có phủ định của $y \in A \land y \in B$. Điều này tương đương với $y \notin A \lor y \notin B$. Đang xét $y \in A$ nên $y \notin B$, dẫn ra được $y \in A \setminus B$.
   \item Nếu $y \in B$, với lập luận tương tự, $y \in B \setminus A$.
\end{itemize}
Do đó, $y \in (A \setminus B) \cup (B \setminus A)$. Vậy,
\begin{equation*}
   (A \cup B) \setminus (A \cap B) \subseteq A \vartriangle B.
\end{equation*}

Từ hai hướng, kết hợp chúng lại để có điều phải chứng minh.

\stepcounter{subexercise}
\arabic{subexercise}. Điều phải chứng minh là khá hiển nhiên, do 
\begin{align*}
   A \vartriangle B &= (A \setminus B) \cup (B \setminus A) \\
      &= (B \setminus A) \cup (A \setminus B) \\
      &= B \vartriangle A.
\end{align*}